%% kapa: exam=1 round=20 year=2009 period=2 form=A subject=accounting subject_ko=회계학 seq=36 number=36
%% source: Q-Net 감정평가사 기출문제 (cst00309) artlSeq=5024178 file=extracted/1008/5024178/2009년 제20회 감정평가사 1차시험 기출문제/2교시.hwp
%% conversion: kordoc 4.24.51 -> tools/gen1.py (auto); review=proofread-sonnet; flags=none
\begin{question}{36}
(주)한국은 제품 A를 생산하며 20×9년 5월초에 영업을 개시하였다(기초재고자산은 없음). 20×9년 5월과 6월의 생산량은 각각 400단위, 500단위이며, 판매량은 각각 380단위, 400단위이다. 매월 고정제조간접원가는 ￦400,000씩 동일하게 발생한다. 20×9년 6월의 전부원가계산에 의한 손익계산서가 다음과 같을 때 6월의 변동원가계산에 의한 영업이익은 얼마인가? (단, 원가흐름가정은 선입선출법을 적용한다.)

{\small\centering
\begin{tabular}{@{}l r r@{}}
\hline
매출액 & & ￦1,000,000 \\
매출원가 & & \\
\hspace{1em}월초제품재고액 & ￦45,000 & \\
\hspace{1em}당월제품제조원가 & 1,050,000 & \\
\hspace{1em}월말제품재고액 & 252,000 & 843,000 \\
\cline{2-2}\cline{3-3}
매출총이익 & & ￦157,000 \\
\hspace{1em}판매비와관리비 & & 67,000 \\
\cline{3-3}
영업이익 & & ￦90,000 \\
\hline
\end{tabular}\par}
\begin{choices}
  \item ￦6,000
  \item ￦14,000
  \item ￦70,000
  \item ￦110,000
  \item ￦166,000
\end{choices}
\end{question}
\answer{2}
